% 乔治·伽莫夫（综述）
% license CCBYSA3
% type Wiki

本文根据 CC-BY-SA 协议转载翻译自维基百科\href{https://en.wikipedia.org/wiki/George_Gamow}{相关文章}

乔治·伽莫夫（George Gamow，有时写作 Gammoff；原名 Гео́ргий Анто́нович Га́мов，Georgiy Antonovich Gamov，1904 年 3 月 4 日－1968 年 8 月 19 日）是一位苏联和美国的博学家、理论物理学家与宇宙学家。他是乔治·勒梅特大爆炸理论的早期倡导者与发展者之一。伽莫夫首先通过量子隧穿理论解释了α衰变的机理，提出了液滴模型——第一个原子核的数学模型，并研究了放射性衰变、恒星形成、恒星核合成、以及大爆炸核合成（他统称为“核宇宙成因论”，nucleocosmogenesis），并预言了宇宙微波背景辐射的存在以及分子遗传学的基础。伽莫夫是量子隧穿现象发展与理解的关键人物之一。

在中晚期生涯中，伽莫夫将大量精力投入教学，并撰写了多部科普著作，包括《一二三……无穷大》以及“汤普金斯先生”系列丛书（1939–1967）。他的一些著作在最初出版半个多世纪后仍在印行。为纪念他，美国科罗拉多大学博尔德分校设立了乔治·伽莫夫纪念讲座。
\subsection{早年生平与职业生涯}
伽莫夫出生于俄国帝国敖德萨（今乌克兰敖德萨）。他的父亲在中学教授俄语与文学，母亲则在女子学校教授地理与历史。除了俄语外，伽莫夫还从母亲那里学会了一些法语，并由家庭教师学习了德语。在大学期间，他又学习了英语，并最终能流利使用。伽莫夫早期的大部分论文以德语或俄语撰写，后来则以英语发表技术论文和科普著作。

他先后就读于敖德萨物理与数学学院（1922–1923），以及列宁格勒大学（1923–1929）。在列宁格勒期间，伽莫夫师从亚历山大·弗里德曼，直到弗里德曼于1925年早逝，被迫更换导师。在大学时期，伽莫夫结识了三位理论物理专业的同学——列夫·朗道、德米特里·伊万年科和马特维·布龙施泰因。四人组成了一个自称为“三剑客”的小组，定期聚会讨论并分析当时发表的开创性量子力学论文。后来，伽莫夫也用同样的称呼来形容他与拉尔夫·阿尔费和罗伯特·赫尔曼共同组成的研究小组。

毕业后，伽莫夫在哥廷根从事量子理论研究，他对原子核结构的研究成果成为其博士论文的基础。此后，他于1928年至1931年在哥本哈根大学理论物理研究所工作，中途曾前往剑桥大学卡文迪许实验室，与著名物理学家欧内斯特·卢瑟福合作。在此期间，伽莫夫继续研究原子核（提出了著名的“液滴模型”），同时还与罗伯特·阿特金森和弗里茨·豪特曼斯合作，研究恒星物理学。

1931年，年仅28岁的伽莫夫被选为苏联科学院通讯院士——成为当时最年轻的院士之一。[5][a][b] 在1931年至1933年间，他在列宁格勒镭研究所物理部工作，该部门由维塔利·赫洛平领导。在伽莫夫、伊戈尔·库尔恰托夫与列夫·米索夫斯基的指导与直接参与下，欧洲第一台回旋加速器的设计工作得以展开。1932年，伽莫夫与米索夫斯基向镭研究所学术委员会提交了加速器设计草案，并获得批准。然而，这台回旋加速器直到1937年才最终建成。[6]
\subsection{放射性衰变}
20世纪早期，科学家已经知道放射性物质具有特征性的指数衰减速率，即所谓的半衰期。与此同时，人们也发现放射性辐射具有特定的能量谱。到1928年时，伽莫夫在哥廷根成功用量子隧穿理论解释了原子核的α衰变现象，并在数学家尼古拉·科钦（的协助下完成了严格的推导。[7][8] 这一问题同时也被罗纳德·W·格尼与爱德华·U·康登独立解决。[9][10] 然而，格尼与康登的结果未能达到伽莫夫那样的定量精度。

在经典物理学中，粒子被束缚在原子核内，是因为要逃离强大的核势阱所需能量极高。按经典观点，必须施加巨大的外力才能将原子核分离，因此这种过程不可能自发发生。
然而，在量子力学框架下，粒子存在一定的概率可以“穿透”势阱的势垒并逃逸。伽莫夫在理论上求解了一个原子核的模型势场，并从第一性原理出发推导出$\alpha$衰变的半衰期与放射能量**之间的定量关系，这一关系此前是通过实验经验总结得到的，即著名的盖革–纳塔尔定律。[11]

若干年后，学界将粒子穿透库仑势垒并发生核反应的概率称为伽莫夫因子或伽莫夫–索末菲因子，以纪念他在量子隧穿理论发展中的奠基性贡献。
\subsection{叛逃}
\begin{figure}[ht]
\centering
\includegraphics[width=8cm]{./figures/f7dbe0721f3a0848.png}
\caption{1931年威廉·布拉格实验室工作人员合影：坐于中央者为威廉·亨利·布拉格，最左为物理学家A. 列别杰夫，最右为乔治·伽莫夫。} \label{fig_QZjmf_1}
\end{figure}
伽莫夫在决定逃离苏联之前，曾在多个苏联科研机构任职。然而，随着政治压迫的日益加剧，他最终选择出逃。1931年，他被正式拒绝前往意大利参加一次科学会议。同年，他与苏联物理学家**柳博芙·沃赫明采娃结婚——伽莫夫亲昵地称她为“Rho”，取自希腊字母 ρ 的发音。此后近两年间，伽莫夫和妻子不断尝试离开苏联，无论是否得到官方许可。在这段时间里，尼尔斯·玻尔及其他朋友多次邀请伽莫夫访问他们所在的国家，但他始终未能获得出境许可。

伽莫夫后来回忆说，他与妻子的前两次叛逃尝试发生在1932年，两次都打算划皮艇出逃：第一次计划横渡黑海前往土耳其，行程约250公里；第二次则打算从摩尔曼斯克出发前往挪威。然而，两次计划都因恶劣天气而失败——幸运的是，他们的举动并未被当局察觉。[12]

1933年，伽莫夫突然被批准前往比利时布鲁塞尔参加第七届索尔维会议。他坚决要求妻子必须同行，甚至表示如果不能与妻子一同前往，他将拒绝单独出席。最终，苏联当局让步，为他们夫妇二人签发了护照。

伽莫夫夫妇顺利参加了会议，并在玛丽·居里及其他物理学家的帮助下，设法延长了在西欧的停留时间。在接下来的一年里，伽莫夫先后在居里研究所、伦敦大学以及密歇根大学获得了临时研究职位。
